%***************************************************************************************************
% preamble/floats-and-tables.tex - podpisy, rysunki, tabele.
% Zasada: podpis tabeli i listingu NAD obiektem, podpis rysunku POD obiektem.
%***************************************************************************************************

% Obrazy wyszukiwane w katalogach images/ i logos/ (dla bezpośredniego \includegraphics)
\graphicspath{{images/}{logos/}}

% Podpisy: mała czcionka, pogrubiona etykieta, kropka ("Rysunek 2.1. Opis")
\captionsetup{font=small, labelfont=bf, labelsep=period, skip=8pt}
\captionsetup[figure]{position=bottom}
\captionsetup[table]{position=top}
\captionsetup[subfigure]{font=footnotesize, labelfont=bf, labelsep=period}

% Domyślne umiejscowienie rysunków (ustawiane w config/metadata.tex)
\floatplacement{figure}{\DefaultFigurePlacement}

% Strony nie są rozciągane do dolnego marginesu - unika "rozjeżdżania" odstępów przy obiektach [H]
\raggedbottom

% Obiekty pływające na górze strony bez dodatkowego odstępu
\makeatletter
\setlength\@fptop{0\p@}
\makeatother

% ---------------------------------------------------------------------------------------------
% Rysunek: \Figure[szerokość]{plik}{podpis}{etykieta}
%   - plik podawany względem katalogu images/ (z rozszerzeniem), np. \Figure{zrzut.png}{Podpis}{fig:zrzut}
%   - szerokość jako ułamek szerokości tekstu (domyślnie 0.9)
%   - gdy pliku brakuje, w dokumencie pojawia się ramka z nazwą pliku (kompilacja nie jest przerywana)
% ---------------------------------------------------------------------------------------------
\newcommand{\Figure}[4][0.9]{%
  \begin{figure}
    \centering
    \IfFileExists{images/#2}{%
      \includegraphics[width=#1\linewidth, height=0.75\textheight, keepaspectratio]{images/#2}%
    }{%
      \PackageWarning{report-template}{Brak pliku obrazu images/#2}%
      \fbox{\parbox[c][5cm][c]{0.8\linewidth}{\centering\textit{Brak pliku obrazu}\\[1ex]\texttt{\detokenize{images/#2}}}}%
    }%
    \caption{#3}\label{#4}%
  \end{figure}%
}

% ---------------------------------------------------------------------------------------------
% Tabele
% ---------------------------------------------------------------------------------------------
\newcommand{\tablefont}{\footnotesize}          % jeden rozmiar czcionki dla wszystkich tabel
\renewcommand{\arraystretch}{1.2}
\setlength{\tabcolsep}{5pt}

% Kolumny o stałej szerokości: L{3cm} (do lewej), C{3cm} (wyśrodkowana), R{3cm} (do prawej)
\newcolumntype{L}[1]{>{\raggedright\arraybackslash}p{#1}}
\newcolumntype{C}[1]{>{\centering\arraybackslash}m{#1}}
\newcolumntype{R}[1]{>{\raggedleft\arraybackslash}p{#1}}
% Kolumny rozciągliwe dla tabularx: Y (do lewej), Z (wyśrodkowana)
\newcolumntype{Y}{>{\raggedright\arraybackslash}X}
\newcolumntype{Z}{>{\centering\arraybackslash}X}
